% Options for packages loaded elsewhere
\PassOptionsToPackage{unicode}{hyperref}
\PassOptionsToPackage{hyphens}{url}
\PassOptionsToPackage{dvipsnames,svgnames,x11names}{xcolor}
\PassOptionsToPackage{space}{xeCJK}
\documentclass[
]{article}
\usepackage{xcolor}
\usepackage[margin=2.5cm]{geometry}
\usepackage{amsmath,amssymb}
\setcounter{secnumdepth}{-\maxdimen} % remove section numbering
\usepackage{iftex}
\ifPDFTeX
  \usepackage[T1]{fontenc}
  \usepackage[utf8]{inputenc}
  \usepackage{textcomp} % provide euro and other symbols
\else % if luatex or xetex
  \usepackage{unicode-math} % this also loads fontspec
  \defaultfontfeatures{Scale=MatchLowercase}
  \defaultfontfeatures[\rmfamily]{Ligatures=TeX,Scale=1}
\fi
\usepackage{lmodern}
\ifPDFTeX\else
  % xetex/luatex font selection
  \setmainfont[]{DejaVu Sans}
  \setmonofont[]{DejaVu Sans Mono}
  \ifXeTeX
    \usepackage{xeCJK}
    \setCJKmainfont[]{Noto Sans CJK SC}
  \fi
  \ifLuaTeX
    \usepackage[]{luatexja-fontspec}
    \setmainjfont[]{Noto Sans CJK SC}
  \fi
\fi
% Use upquote if available, for straight quotes in verbatim environments
\IfFileExists{upquote.sty}{\usepackage{upquote}}{}
\IfFileExists{microtype.sty}{% use microtype if available
  \usepackage[]{microtype}
  \UseMicrotypeSet[protrusion]{basicmath} % disable protrusion for tt fonts
}{}
\makeatletter
\@ifundefined{KOMAClassName}{% if non-KOMA class
  \IfFileExists{parskip.sty}{%
    \usepackage{parskip}
  }{% else
    \setlength{\parindent}{0pt}
    \setlength{\parskip}{6pt plus 2pt minus 1pt}}
}{% if KOMA class
  \KOMAoptions{parskip=half}}
\makeatother
\usepackage{longtable,booktabs,array}
\usepackage{caption}
\captionsetup[table]{skip=6pt}
\newcounter{none} % for unnumbered tables
\usepackage{calc} % for calculating minipage widths
% Correct order of tables after \paragraph or \subparagraph
\usepackage{etoolbox}
\makeatletter
\patchcmd\longtable{\par}{\if@noskipsec\mbox{}\fi\par}{}{}
\makeatother
% Allow footnotes in longtable head/foot
\IfFileExists{footnotehyper.sty}{\usepackage{footnotehyper}}{\usepackage{footnote}}
\makesavenoteenv{longtable}
\setlength{\emergencystretch}{3em} % prevent overfull lines
\providecommand{\tightlist}{%
  \setlength{\itemsep}{0pt}\setlength{\parskip}{0pt}}
\usepackage{bookmark}
\IfFileExists{xurl.sty}{\usepackage{xurl}}{} % add URL line breaks if available
\urlstyle{same}
% fallback for those not using the hyperref driver hyperxmp:
\makeatletter
\@ifundefined{xmpquote}{\newcommand{\xmpquote}[1]{#1}}{}
\makeatother
\hypersetup{
  colorlinks=true,
  linkcolor={Maroon},
  filecolor={Maroon},
  citecolor={Blue},
  urlcolor={Blue},
  pdfcreator={LaTeX via pandoc}}

\author{}
\date{}

\begin{document}

\section{筑基篇课后习题 VI：互锁增长步长与脱离投影------物理空间 Pat
形式化探索}\label{ux7b51ux57faux7bc7ux8bfeux540eux4e60ux9898-viux4e92ux9501ux589eux957fux6b65ux957fux4e0eux8131ux79bbux6295ux5f71ux7269ux7406ux7a7aux95f4-pat-ux5f62ux5f0fux5316ux63a2ux7d22}

\begin{quote}
\textbf{声明 (Declaration)}: 本文\textbf{不代表任何数学结论}
(non-mathematical-claim), 仅为基于 Lean 4
形式化的一种\textbf{实验观测报告} (experimental observation report)。

{\def\LTcaptype{none} % do not increment counter
\begin{longtable}[]{@{}
  >{\raggedright\arraybackslash}p{(\linewidth - 6\tabcolsep) * \real{0.2500}}
  >{\raggedright\arraybackslash}p{(\linewidth - 6\tabcolsep) * \real{0.2500}}
  >{\raggedright\arraybackslash}p{(\linewidth - 6\tabcolsep) * \real{0.2500}}
  >{\raggedright\arraybackslash}p{(\linewidth - 6\tabcolsep) * \real{0.2500}}@{}}
\toprule\noalign{}
\begin{minipage}[b]{\linewidth}\raggedright
习题
\end{minipage} & \begin{minipage}[b]{\linewidth}\raggedright
主题
\end{minipage} & \begin{minipage}[b]{\linewidth}\raggedright
Lean 形式化
\end{minipage} & \begin{minipage}[b]{\linewidth}\raggedright
DOI
\end{minipage} \\
\midrule\noalign{}
\endhead
\bottomrule\noalign{}
\endlastfoot
exercise\_i\_6 & 互锁增长步长与脱离投影（R161/R162/R163） &
PatInterlockGrowth.lean / PatThreeBodyLocal.lean /
PatThreeBodyShared.lean & 10.5281/zenodo.21916850 \\
\end{longtable}
}
\end{quote}

\textbf{Exercise VI for the Foundation-Building Chapter: Interlock
Growth Steps and Detachment Projection --- PAT Formalization of Physical
Space}

\begin{quote}
习题定位：筑基篇课后习题系列第六题（用户两个问题的直接探索）。① 合法步长
6 是否真的存在------怀疑是 2
的倍数，与物理世界''未经基点还原的对称性增长''的关系；②
三互锁基础上，一对互锁彻底脱离物理空间后，剩余部分遵守物理空间法则的情况------力学/量子力学/引力的
Pat
形式化纲领。全部新定理只锚筑基篇定理（R136/R138/R143/R144/R149/R160），不用外部引理。
\end{quote}

\emph{2026-08-13 · Internal research exercise · Lean 4 / mathlib v4.32.2
· 7 new theorems PROVED no sorry · 算器神魂论筑基篇课后习题 VI }

\begin{center}\rule{0.5\linewidth}{0.5pt}\end{center}

\subsection{习题陈述}\label{ux4e60ux9898ux9648ux8ff0}

\begin{quote}
① 探索合法步长：6 互锁是否真的存在？怀疑可能是 2
的倍数------这与物理世界下''未经基点还原的对称性增长''有什么关系？②
观测：三互锁基础上，若有一对互锁彻底脱离物理空间，剩下三对互锁在物理空间投影结构丢失后，剩余部分遵守物理空间法则的情况------特别是物理空间力学、量子力学、引力的
Pat 形式化。
\end{quote}

\textbf{唯一论点}：① \textbf{合法步长 = 2 的倍数（2, 4, 6, 8,
\ldots），不是 2 的幂------6
互锁存在}，因为三体断裂是闭合回路共享端点的特殊病（线性相关，R160），不是奇数互锁本身不可行；3
对独立正交轴的互锁（θ₁, θ₂, θ₃ 互不约束）全部自洽。②
\textbf{互锁逐对独立}------任何一对的脱离不影响其余对的互锁性；3 对脱离
1 对 = 剩 2 对（4 互锁 = S³），物理空间剩余结构遵守物理空间法则。

\begin{center}\rule{0.5\linewidth}{0.5pt}\end{center}

\subsection{零、总论：三体断裂不是''奇数病''，是''闭合回路病''}\label{ux96f6ux603bux8bbaux4e09ux4f53ux65adux88c2ux4e0dux662fux5947ux6570ux75c5ux662fux95edux5408ux56deux8defux75c5}

R160 证明了三互锁闭合回路断裂：三点共享端点 ⟹ 相位差之和恒为 0 ⟹
三方向线性相关 ⟹ 无法独立锁定。\textbf{但关键区分}：

\begin{itemize}
\tightlist
\item
  \textbf{闭合回路病}：3 个互锁对共享 3 个端点 (e₁, e₂, e₃)，d₁₂ + d₂₃ +
  d₃₁ = 0 恒成立（R160）------这是共享端点的特殊结构病。
\item
  \textbf{独立轴互锁}：3 对互锁作用在独立正交轴上，θ₁, θ₂, θ₃
  互不约束------每对独立成对 exp(iθⱼ)·exp(-iθⱼ) =
  1，没有闭合回路线性相关。
\end{itemize}

⟹ \textbf{6 互锁（3 对独立轴）存在}。合法步长 = 2 的倍数，不是 2
的幂。三体不可解（Poincaré）的结构对应 =
闭合回路（共享端点）的线性相关，不是''3 互锁''本身。

\begin{center}\rule{0.5\linewidth}{0.5pt}\end{center}

\subsection{一、★6 互锁存在：3 对独立相位互锁（合法步长 = 2
的倍数）}\label{ux4e006-ux4e92ux9501ux5b58ux57283-ux5bf9ux72ecux7acbux76f8ux4f4dux4e92ux9501ux5408ux6cd5ux6b65ux957f-2-ux7684ux500dux6570}

\textbf{★核心：任意 k 对独立相位互锁全部自洽------合法步长是 2
的倍数（2, 4, 6, 8, \ldots），不是 2 的幂；6 互锁（3 对）存在。}

\subsubsection{论证}\label{ux8bbaux8bc1}

\begin{enumerate}
\def\labelenumi{\arabic{enumi}.}
\tightlist
\item
  \textbf{互锁成对}（R136 ②③）：互锁的最小单位 = 一对对称性 \{d,
  -d\}。每次增长 +1 对 = +2 相位方向。
\item
  \textbf{单对自洽}（R138）：exp(iθ)·exp(-iθ) = exp(0) =
  1------相位差可加（RulerPhase：相位差 = 方向），对称对还原到
  1（R143）。
\item
  \textbf{独立轴无闭合回路}：3 对独立相位 (θ₁, θ₂, θ₃)
  互不约束------每对 exp(iθⱼ)·exp(-iθⱼ) = 1 只依赖自己的
  θⱼ，没有三体闭合回路的线性相关（对照 R160）。
\item
  \textbf{任意 k 对}：∀ k, ∀ θ : Fin k → ℝ, ∀ j, exp(iθⱼ)·exp(-iθⱼ) =
  1------合法步长 = 2 的倍数（2, 4, 6, 8, \ldots），\textbf{6
  互锁存在}。
\item
  \textbf{与物理对称性增长的关系}（未经基点还原）：互锁数 =
  成对的对称性方向数。未经基点还原的对称性可任意成对增长（2
  的倍数）------每次增加一对对称方向（+2）；还原后坍缩到折叠类（R144：对称对还原到锚点
  0/1）。\textbf{物理世界中''未经基点还原的对称性增长'' =
  互锁对的成对添加}------这解释了为什么对称性方向总是成对出现（电荷共轭/宇称/时间反演
  CPT 对称、粒子-反粒子对等）。
\end{enumerate}

\subsubsection{形式化（PatInterlockGrowth.lean，4
定理）}\label{ux5f62ux5f0fux5316patinterlockgrowth.lean4-ux5b9aux7406}

\begin{itemize}
\tightlist
\item
  \texttt{pair\_interlock\_self\_consistent}：exp(iθ)·exp(-iθ) =
  1------单对互锁自洽。
\item
  \texttt{three\_independent\_pairs\_interlock}：★6 互锁存在------3
  对独立相位互锁全部自洽。
\item
  \texttt{k\_pairs\_independent\_interlock}：★合法步长 = 2
  的倍数------任意 k 对独立互锁自洽。
\item
  \texttt{three\_body\_loop\_special}：三体闭合回路是特殊病（线性相关），独立轴互锁不受此限。
\end{itemize}

\textbf{验证}：0 sorry。单对自洽用 exp\_add + exp\_zero；k 对用 Fin
量化。

\begin{center}\rule{0.5\linewidth}{0.5pt}\end{center}

\subsection{二、★脱离投影：一对互锁脱离物理空间，剩余对仍互锁}\label{ux4e8cux8131ux79bbux6295ux5f71ux4e00ux5bf9ux4e92ux9501ux8131ux79bbux7269ux7406ux7a7aux95f4ux5269ux4f59ux5bf9ux4ecdux4e92ux9501}

\textbf{★核心：互锁逐对独立------任何一对的脱离不影响其余对的互锁性；3
对脱离 1 对 = 剩 2 对（4 互锁 =
S³），物理空间剩余结构遵守物理空间法则。}

\subsubsection{论证}\label{ux8bbaux8bc1-1}

\begin{enumerate}
\def\labelenumi{\arabic{enumi}.}
\tightlist
\item
  \textbf{逐对独立性}：每对互锁 exp(iθⱼ)·exp(-iθⱼ) = 1 只依赖自己的
  θⱼ------\textbf{一对的脱离不影响其他对的互锁性}。
\item
  \textbf{脱离投影}：3 对 (θ₁, θ₂, θ₃) 中第 3 对彻底脱离物理空间（θ₃
  进入高维/不可观测方向），物理空间剩 (θ₁, θ₂)
  两对------仍满足互锁（第一节）。
\item
  \textbf{剩余结构 = 4 互锁 = S³}：剩 2 对 = R149 四相位（2 对）→ 归一化
  S³（R154）------物理空间剩余结构遵守物理空间法则（互锁保持性）。
\item
  \textbf{一般化}：任意 k 对中脱离 1 对，剩余 k-1
  对仍互锁------\textbf{投影（脱离某些对）保持剩余对的互锁}，这是''投影结构丢失后剩余部分遵守物理空间法则''的形式化。
\item
  \textbf{物理图景}：3 对（6 互锁 = 三维空间 3 轴？）中 1
  对脱离（进入高维）⟹ 物理空间剩 2 对 = 4 互锁 = S³ = SU(2)
  结构------\textbf{物理空间的量子结构（S³/SU(2)）可能是 6 互锁（3
  对）投影掉 1 对后的剩余}。
\end{enumerate}

\subsubsection{形式化（PatInterlockGrowth.lean，2
定理）}\label{ux5f62ux5f0fux5316patinterlockgrowth.lean2-ux5b9aux7406}

\begin{itemize}
\tightlist
\item
  \texttt{pair\_detachment\_remaining\_interlocked}：★脱离投影（3 对中 1
  对脱离）------剩余两对仍互锁（4 互锁 = S³）。
\item
  \texttt{pair\_detachment\_general}：★脱离投影一般化------任意 k
  对中脱离某些对，剩余对仍互锁（逐对独立）。
\end{itemize}

\textbf{验证}：0 sorry。逐对独立性直接由单对自洽定理给出。

\begin{center}\rule{0.5\linewidth}{0.5pt}\end{center}

\subsection{三、物理空间力学/量子力学/引力的 Pat
形式化纲领（CONJECTURE/OBSERVATION
标注）}\label{ux4e09ux7269ux7406ux7a7aux95f4ux529bux5b66ux91cfux5b50ux529bux5b66ux5f15ux529bux7684-pat-ux5f62ux5f0fux5316ux7eb2ux9886conjectureobservation-ux6807ux6ce8}

\begin{quote}
本节为纲领性探索（非定理），标注认识论标签：已形式化部分 =
OBSERVATION（锚定定理），未形式化部分 = CONJECTURE（待后续习题）。
\end{quote}

\subsubsection{3.1 力学（Kepler 二体 =
单互锁对，已覆盖）}\label{ux529bux5b66kepler-ux4e8cux4f53-ux5355ux4e92ux9501ux5bf9ux5df2ux8986ux76d6}

\begin{itemize}
\tightlist
\item
  \textbf{二体可解} =
  单互锁对（R147：互锁成对；R138：相位锁定）------Kepler 椭圆轨道 =
  相位锁定的单互锁对。
\item
  \textbf{三体不可解} = 闭合回路三互锁线性相关（R160；三体习题）。
\item
  \textbf{脱离投影的力学图景}：3 体中 1 对脱离（两个天体不再相互作用）⟹
  剩 2 对------但闭合回路病仍在（共享端点）。力学意义的脱离 =
  天体对解耦（第三体远离），此时闭合回路退化为两条独立路径------\textbf{逐对独立性使解耦后的天体遵守独立二体法则}（OBSERVATION：与
  pair\_detachment\_general 对应）。
\end{itemize}

\subsubsection{3.2 量子力学（S³ = SU(2)
自旋结构，纲领）}\label{ux91cfux5b50ux529bux5b66suxb3-su2-ux81eaux65cbux7ed3ux6784ux7eb2ux9886}

\begin{itemize}
\tightlist
\item
  \textbf{4 互锁 = S³ = SU(2)}（R149/R154 已证）：4 相位互锁归一化 =
  单位 3 维球面------\textbf{SU(2)（自旋群）的 Pat 结构对应 = 4
  互锁}（OBSERVATION：锚定 R149/R154）。
\item
  \textbf{6 互锁 = 3 对 = 三维空间 3 轴}（本习题）：3
  对独立轴互锁------\textbf{三维物理空间的 Pat 结构对应 = 6
  互锁}（CONJECTURE：三维性来源）。
\item
  \textbf{量子 = 6 投影掉 1
  对}（CONJECTURE）：物理空间的量子结构（S³/SU(2)，4 互锁）可能是 6
  互锁（三维空间）投影掉 1
  对后的剩余------\textbf{自旋-空间对应（spin-space correspondence）的
  Pat 表达}：SU(2) 是 SO(3) 的双重覆盖，4 互锁（SU(2)）是 6 互锁（3
  对）脱离 1 对后的结构。
\item
  \textbf{泡利矩阵/四元数}（CONJECTURE）：4 互锁的 2 对 = 1 轴（数值）+
  i 轴（相位）------四元数 \{1, i, j, k\} 的 Pat 对应，j, k
  是额外脱离的对。
\end{itemize}

\subsubsection{3.3 引力（脱离对 =
高维方向，纲领）}\label{ux5f15ux529bux8131ux79bbux5bf9-ux9ad8ux7ef4ux65b9ux5411ux7eb2ux9886}

\begin{itemize}
\tightlist
\item
  \textbf{引力 = 最大脱离}（CONJECTURE）：一对互锁彻底脱离物理空间 =
  高维方向（超出物理空间可观测维度）------引力（时空曲率）的 Pat 对应 =
  脱离对与剩余对的耦合方式。
\item
  \textbf{投影保持互锁}（OBSERVATION：本习题
  pair\_detachment\_general）：脱离后剩余部分仍遵守物理空间法则------\textbf{物理定律在投影下的保持性}是
  Pat 的逐对独立性。
\item
  \textbf{多维时空}（CONJECTURE）：3 空间对 + 1 时间对 = 4 对 = 8
  互锁？------时空维数的 Pat 计数。
\end{itemize}

\subsubsection{诚实边界}\label{ux8bdaux5b9eux8fb9ux754c}

\begin{itemize}
\tightlist
\item
  已形式化：步长 2 的倍数（k
  对独立互锁自洽，PROVED）、脱离投影保持（逐对独立，PROVED）、三体闭合回路特殊（R160）。
\item
  未形式化（CONJECTURE）：三维空间来源、自旋-空间对应、引力脱离耦合、时空维数------这些需要新的结构定义（三维正交轴、SU(2)
  表示、时空度量），不在本习题能力内，标注为后续习题方向。
\end{itemize}

\begin{center}\rule{0.5\linewidth}{0.5pt}\end{center}

\subsection{四、习题解答总结}\label{ux56dbux4e60ux9898ux89e3ux7b54ux603bux7ed3}

{\def\LTcaptype{none} % do not increment counter
\begin{longtable}[]{@{}
  >{\raggedright\arraybackslash}p{(\linewidth - 6\tabcolsep) * \real{0.2500}}
  >{\raggedright\arraybackslash}p{(\linewidth - 6\tabcolsep) * \real{0.2500}}
  >{\raggedright\arraybackslash}p{(\linewidth - 6\tabcolsep) * \real{0.2500}}
  >{\raggedright\arraybackslash}p{(\linewidth - 6\tabcolsep) * \real{0.2500}}@{}}
\toprule\noalign{}
\begin{minipage}[b]{\linewidth}\raggedright
论点
\end{minipage} & \begin{minipage}[b]{\linewidth}\raggedright
内容
\end{minipage} & \begin{minipage}[b]{\linewidth}\raggedright
锚到
\end{minipage} & \begin{minipage}[b]{\linewidth}\raggedright
定理
\end{minipage} \\
\midrule\noalign{}
\endhead
\bottomrule\noalign{}
\endlastfoot
★6 互锁存在 & 3 对独立相位互锁全部自洽------合法步长 = 2 的倍数，非 2
的幂 & R136/R138/R143 & three\_independent\_pairs\_interlock /
k\_pairs\_independent\_interlock \\
三体断裂 = 特殊病 & 闭合回路共享端点 ⟹ 线性相关；独立轴互锁不受此限 &
R160 & three\_body\_loop\_special \\
★脱离投影 & 3 对中 1 对脱离 ⟹ 剩 2 对仍互锁（4 互锁 = S³）；任意 k
对逐对独立 & R149/R154 & pair\_detachment\_remaining\_interlocked /
pair\_detachment\_general \\
物理纲领 & 力学（解耦=逐对独立）/量子（S³=SU(2) 4
互锁）/引力（脱离对=高维） & R149/R154 & 3.1--3.3（CONJECTURE 标注） \\
\end{longtable}
}

\textbf{习题的回答（两句话）}： 1. \textbf{合法步长 = 2 的倍数，6
互锁存在}------三体断裂是闭合回路共享端点的特殊病（线性相关，R160），不是奇数互锁本身不可行；任意
k 对独立轴互锁全部自洽。未经基点还原的对称性增长 =
互锁对的成对添加（每次 +1 对 = +2 方向）。 2.
\textbf{互锁逐对独立}------任何一对脱离物理空间不影响其余对的互锁性；3
对脱离 1 对 = 剩 2 对（4 互锁 =
S³），物理空间剩余结构遵守物理空间法则。物理空间的量子结构（S³/SU(2)）可能是
6 互锁（三维）投影掉 1 对后的剩余（CONJECTURE）。

\begin{center}\rule{0.5\linewidth}{0.5pt}\end{center}

\subsection{五、相关对照}\label{ux4e94ux76f8ux5173ux5bf9ux7167}

{\def\LTcaptype{none} % do not increment counter
\begin{longtable}[]{@{}
  >{\raggedright\arraybackslash}p{(\linewidth - 4\tabcolsep) * \real{0.3333}}
  >{\raggedright\arraybackslash}p{(\linewidth - 4\tabcolsep) * \real{0.3333}}
  >{\raggedright\arraybackslash}p{(\linewidth - 4\tabcolsep) * \real{0.3333}}@{}}
\toprule\noalign{}
\begin{minipage}[b]{\linewidth}\raggedright
文献
\end{minipage} & \begin{minipage}[b]{\linewidth}\raggedright
对应内容
\end{minipage} & \begin{minipage}[b]{\linewidth}\raggedright
备注
\end{minipage} \\
\midrule\noalign{}
\endhead
\bottomrule\noalign{}
\endlastfoot
\textbf{Poincaré, H.} 1890. \emph{Sur le problème des trois corps} &
三体不可积性 & 闭合回路病的经典对应 \\
\textbf{Pauli, W.} 1927. \emph{Zur Quantenmechanik des magnetischen
Elektrons} & 自旋/泡利矩阵 & S³ = SU(2) 的量子对应（CONJECTURE 层） \\
\textbf{Hurwitz, A.} 1898. \emph{Über die Composition der quadratischen
Formen} & 可除代数 ℝ, ℂ, ℍ, 𝕆（1,2,4,8 维） & 步长 2
的幂的经典对应（对照：本习题证明步长是 2 的倍数更宽） \\
\textbf{R136/R138/R143/R144/R149/R160}（筑基篇） &
互锁成对/相位锁定/对称对还原/四相位/三体断裂 & 本框架定理，非外部引理 \\
\end{longtable}
}

\begin{center}\rule{0.5\linewidth}{0.5pt}\end{center}

\emph{筑基篇课后习题 VI · 2026-08-13 · Lean 4 / mathlib v4.32.2 · 7
theorems PROVED no sorry（PatInterlockGrowth.lean）· 配套 Lean
形式化：formal/Formal/Toolkit/PatInterlockGrowth.lean（快照见
../lean/PatInterlockGrowth.lean）}

\begin{center}\rule{0.5\linewidth}{0.5pt}\end{center}

\subsection{六、★互锁是本地操作：3 组互锁的最小锁定数量 =
12（用户修正）}\label{ux516dux4e92ux9501ux662fux672cux5730ux64cdux4f5c3-ux7ec4ux4e92ux9501ux7684ux6700ux5c0fux9501ux5b9aux6570ux91cf-12ux7528ux6237ux4feeux6b63}

\textbf{★核心：互锁是本地的------每组的互锁只依赖自己的 (aⱼ,
θⱼ)，不与其他组共享；每组最少 4 互锁（否则被 R134 吸收）；3
组互锁的最小锁定数量 = 3 × 4 = 12。}

\subsubsection{论证}\label{ux8bbaux8bc1-2}

\begin{enumerate}
\def\labelenumi{\arabic{enumi}.}
\tightlist
\item
  \textbf{互锁是本地操作}（用户洞察修正）：每组的互锁只依赖自己的参数------R149
  quadriphase\_interlock 逐组版：∀ j, aⱼ·(1/aⱼ) = 1 ∧ exp(iθⱼ)·exp(-iθⱼ)
  = 1 ∧ log 对 ∧ 范数对。互锁对是成对的局部操作（R136
  ②③：方向成对声明只涉及自己的方向），\textbf{各组互不依赖，不共享端点}。
\item
  \textbf{每组最少 4 互锁}（R160）：4 互锁是最小自洽互锁结构（1 退化
  R062，3 断裂 R160）------\textless{} 4 的组会被 pat0 吸收（R134：app
  pat0 pat0 = pat0, layerUp pat0 = pat0）。
\item
  \textbf{★3 组互锁的最小锁定数量 = 3 × 4 = 12}：3 组本地互锁 × 每组最少
  4 = 12 互锁；12 = 6 对（成对性满足）；12 互锁自洽（R161：任意 k
  对独立互锁自洽，12 = 6 对）。
\item
  \textbf{与全局计数的区分}：R160 的''三互锁断裂''是 3
  个互锁对\textbf{共享 3 端点}的全局闭合回路（线性相关）；本地视角下 3
  组互锁 = 3 个独立的 4 互锁组（每组自己的 (aⱼ,
  θⱼ)），\textbf{不共享端点------无闭合回路病}。
\item
  \textbf{三体的本地图景}：三体 = 3 个本地 4 互锁单体（S³，各 4 互锁 =
  12
  总互锁）------存在性由本地互锁保证（不被吸收）；三体不可解（Poincaré）是组间相互作用的动力学非线性，不是单体互锁断裂。
\end{enumerate}

\subsubsection{形式化（PatThreeBodyLocal.lean，5
定理）}\label{ux5f62ux5f0fux5316patthreebodylocal.lean5-ux5b9aux7406}

\begin{itemize}
\tightlist
\item
  \texttt{interlock\_is\_local}：★互锁是本地操作------任意 k
  组互锁逐组独立 4 互锁（R149 逐组版）。
\item
  \texttt{body\_minimal\_four}：每组最少 4 互锁（R160
  最小自洽，\textless{} 4 被 R134 吸收）。
\item
  \texttt{three\_groups\_minimal\_twelve}：★3 组互锁的最小锁定数量 =
  12（3 × 4）。
\item
  \texttt{twelve\_interlock\_self\_consistent}：12 互锁自洽（6
  对独立，R161）。
\item
  \texttt{three\_body\_local\_perspective}：★全景------互锁本地 ∧
  每组最少 4 ∧ 3 组 = 12。
\end{itemize}

\textbf{验证}：0 sorry。锚定 R149 / R160 / R161 / R134。

\subsubsection{连接}\label{ux8fdeux63a5}

这修正了习题 V（R160）的视角：三互锁断裂不是''3
互锁''本身不可行，而是\textbf{全局共享端点的闭合回路病}。本地视角下，每个单体独立
4 互锁（S³），三体 = 12 互锁（6
对），自洽存在------\textbf{三体不可解（Poincaré）是组间相互作用的动力学非线性，与单体互锁结构无关}。

\begin{center}\rule{0.5\linewidth}{0.5pt}\end{center}

\emph{筑基篇课后习题 VI（增补）· 2026-08-13 · 配套 Lean
形式化：formal/Formal/Toolkit/PatThreeBodyLocal.lean（快照见
../lean/PatThreeBodyLocal.lean）}

\end{document}
